El repositorio de la tesis conserva una separación explícita entre el manuscrito y los artefactos que lo sustentan. La estructura inicial es:

\begin{verbatim}
Master-s-thesis/
  thesis/              # Fuente LaTeX y PDF generado
  notebooks/           # Exploración, validación y experimentos
  data/raw/            # Datos originales; no se versionan datos restringidos
  data/processed/      # Transformaciones trazables
  src/                 # Código reusable que emerge de los notebooks
  docs/                # Notas metodológicas y de reproducibilidad
  apps/web/            # Sitio público y visor del artefacto
  VERSION              # Versión común de los artefactos
  VERSION_HISTORY.json # Registro de cambios
\end{verbatim}

Para construir el manuscrito se ejecuta \texttt{make pdf} desde el directorio \texttt{thesis/}. La integración de GitHub Pages compila el documento y publica el PDF junto con documentación y notebooks. Los datos, las credenciales y los resultados sensibles permanecen fuera del repositorio público.
